\documentclass[10pt]{article}

\usepackage[utf8]{inputenc}

\usepackage[T1]{fontenc}

\usepackage{amsmath}

\usepackage{amsfonts}

\usepackage{amssymb}

\usepackage[version=4]{mhchem}

\usepackage{stmaryrd}



\title{0 }



\author{}

\date{}





\begin{document}

\maketitle

ОБЛАСТИ СО СЛОХНОЙ СТРУКТУРОЙ - области (среды), либо состоящие из большого числа мелких непересекающихся компонентов с различными свойствами, либо имеющие границу сложной (мелкозернистой) структуры (состояшую из большого числа мелких компонентов). Характерной особенностью таких сред является то, что размеры неоднородностей (компонентов) весьма малы по сравнению с размерами всей области, но велики по сравнению с межмолекулярными расстояниями, что позволяет описывать вещество внутри таких неоднородностей уравнениями физики сплошных сред. Необходимость решения краевых задач для уравнений математич. физики в областях со сложной структурой возникает в целом ряде задач физики, механики, химии и биологии. Примерами таких областей могут служить гетерогенные и пористые среды, композиционные материалы, мембраны различных типов [напр., применяемые для разделения жидких (газовых) смесей или биологические].



ОБЛАСти сО слОхной СТрУКтУРОЙ; крае вые задачи для математической физики-область математических исследований, имеющая обширные приложения к различным задачам физики, механики, химии и биологии. Физич. процессы в сильно неоднородных оредах описываются уравнениями с частными производными с быстро изменяющимися (осциллирующими) коэффициентами либо решениями краевых задач для таких уравнений в нек-рой области с граничными условиями, заданными на границе, к-рая имеет мелкозернистую структуру. Предполагается, что характерные размеры неоднородностей имеют порядок $\varepsilon$ такой, что $\varepsilon \ll L$, где $L-$ характерный размер рассматриваемой области $\Lambda$, а сами неоднородности размещены в $\Lambda$ в соответствии с нек-рым законом (правилом) периодически, квазипериодически, случайно (с нек-рой функщией распределения).



Сложная структура области обычно не вносит дополнительных трудностей в доказательства теорем существования и единственности решений таких краевых задач, но, в то же время, нахождение аналитических (и даже численных) точных или приближенных решений таких задач сопряжено с трудностями. Наиболее распространен следующий метод исследования таких задач.



Изучается поведение решений исходной краевой задачи при $\varepsilon \rightarrow 0$ и строится так наз. усредненное уравнение (или краевая задача), к-рым удовлетворяют предельные функции. При этом в одних случаях решения исходной краевой задачи заменяются решениями усредненных (то есть более простых - часто с постоянными коэффициентами) дифференциальных уравнений, рассматриваемых во всем пространстве, в других - исходная сложная граница заменяется средней простой поверхностью, на к-рой задаются так наз. усредненные граничные условия. Таким образом, исходная задача сводится фактически к изучению однородных сред (или сред с однородной границей), описываемых так наз. эффективными характеристиками сред со сложной структурой, для нахождения к-рых существуют хорошо разработанные как математические, так и численные методы. Изложенный метод решения краевых задач в О. со с. с. является



394 ОБЛАСТИ универсальным и применим к различным процессам, протекающим в средах со сложной структурой: к упругим колебаниям, теплопроводности, диффузии, фильтрации, к задачам электростатики, к распространению электромагнитных волн (напр., к задаче дифракции), излучению и т. п. При этом могут рассматриваться как линейные, так и нелинейные модели, как дифференциальные, так и операторные уравнения, а также краевые задачи различных типов (Дирихле, Неймана и др.). Данный метод может быть строго обоснован и позволяет в ряде случаев оценивать погрешность полученного решения.



В ряде случаев (напр., в задачах электростатики, диффузии, фильтрации и др.) можно построить и исследовать решения исходной краевой задачи в виде разложения по степеням малого параметра $\varepsilon / L_{0}\left(L_{0}\right.$ - характерное расстояние между неоднородностями), при этом можно, в принципе, находить такое решение с точностью до $\left(\varepsilon / L_{0}\right)^{n}$, где $n$ - произвольное целое число, и учитывать эффекты, связанные с микроструктурой исследуемой области.



Исследование краевых задач в средах со сложной структурой восходит к работам О. Моссотти (О. Mossotti), Дж. Максвелла (J. Maxwell), Дж. Рэлея (J. Rayleigh).



Jum.: [1] Maxwell J. C., Treatise on Electricity and Magnetismus, v. 1, Oxford, 1873; [2] Rayleigh J.W., «Philos. Mag.», 1892, v. 34, p. 481; [3] Марченко В.А., Хруслов Е.Я., Краевые задачи в областях с мелкозернистой границей, Киев, 1974; [4] Санче сПаленсия Э., Неоднородные среды и теория колебаний, пер. с англ., М., 1984; [5] Бахвалов Н. С., Панасенко Г. П., Осреднение процессов в периодических средах, М., 1984. П. В. Малышев. ОБЛАСТИ СО СЛОЖНОЙ СТРУКТУРОЙ; э ф Ф е к т и в ны е характе р и стики - характеристики, задача расчета которых возникает при исследовании краевых задач для уравнений математической физики в областях со сложной структурой, когда решения исходной краевой задачи в среде со сложной структурой (состоящей из смеси частиц с различными свойствами) заменяются решениями модельной (более простой) задачи в однородной среде, характеризующейся некоторыми эффективными коэффициентами. Причем эффективные характеристики сред со сложной структурой необходимо подбирать так, чтобы решения модельной задачи были близки к решениям исходной. При расчете эффективных характеристик используются как аналитические, так и численные методы.



Ниже приведены формулы для определения эффективной диэлектрич. проницаемости среды, состоящей из смеси однородных частиц во внешнем электрич. поле (аналогичные формулы можно получить для магнитной проницаемости среды, коэффициентов диффузии и фильтрации; используемый метод является универсальным, то есть может быть применен и в ряде других задач). Эта задача восходит к середине и 2-й половине 19 в. к работам О. Моссотти (O. Mossotti), Р. Клаузиуса (R. Clausius), Дж. Максвелла (J. Maxwell), Г. Лоренца (H. Lorentz), Дж. Рэлея (J. Rayleigh).



Для системы сферич. тел с диэлектрич. проницаемостью $\varepsilon_{2}$, расположенных в среде с диэлектрич. проницаемостью $\varepsilon_{1}$ периодически или случайным образом, Дж. Рэлеем (1892) была установлена (с помощью разложения решений





\end{document}